% ============================== FRONT MATTER =========================
\chapter{Before You Start}

This booklet takes the AMATH301 mock exam (three problems, two hours, 20 points) and goes through
it one question at a time. For every question you get the statement word for word, the idea from
the course that the question is really about, how to get started when you are sitting in front of
it, the answer written the way you would write it on your copy, and the mistakes that cost marks.
Every number in it was checked by computer.

You do not need anything else open while you read. When a question uses a theorem, the theorem is
stated right there, with its conditions.

\section{How each question is laid out}

{\renewcommand{\thequ}{\unskip}
\begin{question}{Problem 1 \;·\; 8 pts}
The exam question, copied exactly as it appears on the paper.
\end{question}}\setcounter{qu}{0}

\begin{lesson}
The idea from the lectures that the question tests, explained from zero. If you only have time
to reread one box per question, reread this one.
\end{lesson}

\begin{howto}
The first two or three things to write down. This is what gets you moving when your head goes blank.
\end{howto}

\begin{answer}
The full answer, written the way you would write it on the copy, with every justification the
marker looks for.
\end{answer}

\begin{trap}
A mistake that looks right and loses marks.
\end{trap}

\begin{intuition}
A picture or an analogy for the same idea, when it helps.
\end{intuition}

\section{What this mock is really testing}

The three problems look different, but each one is built around a single lesson of the course.

\begin{center}
\begin{tikzpicture}[
   lec/.style={draw=tealD, fill=ice!30, rounded corners=2mm, text width=3.6cm, align=left,
     font=\sffamily\scriptsize, inner sep=5pt, line width=0.7pt},
   pb/.style={draw=coralD, fill=blush!28, rounded corners=2mm, text width=4.6cm, align=left,
     font=\sffamily\scriptsize, inner sep=5pt, line width=0.7pt},
   arr/.style={-{Stealth[length=2mm]}, tealD, line width=0.8pt}]
  \node[lec] (L1) at (0,0) {\textbf{\color{tealDD}Lecture I}\\ domains, open and closed sets,
     level sets, partial derivatives, differentiability, $C^1$};
  \node[lec] (L2) at (0,-2.4) {\textbf{\color{tealDD}Lecture II}\\ chain rule, tangent planes,
     directional derivatives, gradient};
  \node[lec] (L3) at (0,-4.6) {\textbf{\color{tealDD}Lecture III}\\ second derivatives, the
     Laplacian (TD 3)};
  \node[pb] (P1) at (6.3,0.3) {\textbf{\color{coralD}Problem 1} \;(8 pts)\\ One function,
     $\sqrt{2y-x^2}$, and every tool of Lectures I and II applied to it. \emph{Lesson: the gradient
     tells you everything near a point.}};
  \node[pb] (P2) at (6.3,-2.4) {\textbf{\color{coralD}Problem 2} \;(7 pts)\\ The classic
     counterexample $\frac{xy^2}{x^2+y^2}$. \emph{Lesson: having derivatives in every direction is not
     the same as being differentiable.}};
  \node[pb] (P3) at (6.3,-4.9) {\textbf{\color{coralD}Problem 3} \;(5 pts)\\ Radial functions
     $\varphi(r)$ and the Laplacian. \emph{Lesson: the chain rule, used twice, turns a PDE into an ODE.}};
  \draw[arr] (L1.east) -- (P1.west);
  \draw[arr] (L2.east) -- (P1.west);
  \draw[arr] (L1.east) -- (P2.west);
  \draw[arr] (L2.east) -- (P2.west);
  \draw[arr] (L2.east) -- (P3.west);
  \draw[arr] (L3.east) -- (P3.west);
\end{tikzpicture}
\end{center}

So the whole paper rests on four things you must be able to write from memory: the definition of
differentiability, the $C^1$ theorem, the chain rule, and the formula $D_uf=\grad f\cdot u$ for a
unit vector $u$. Each of them is stated in full when it first comes up.

\section{Your two hours}

Twenty points in 120 minutes is six minutes per point. That gives a plan you can decide now, so
you do not have to decide it in the room.

\begin{center}
\begin{tikzpicture}[x=0.105cm]
  \foreach \a/\b/\c/\t in {0/5/ice/Read,5/50/aqua/Problem 1 \;·\; 45 min,
                           50/90/sky/Problem 2 \;·\; 40 min,90/112/coral/Problem 3 \;·\; 22 min,
                           112/120/blush/Check}{
    \fill[\c, rounded corners=1pt] (\a,0) rectangle (\b,0.8);
    \node[font=\sffamily\bfseries\tiny, text=tealDD, align=center] at ({(\a+\b)/2},0.4) {\t};}
  \foreach \m in {0,30,60,90,120}{\draw[tealDD] (\m,0) -- (\m,-0.15)
    node[below, font=\sffamily\tiny, text=tealDD] {\m'};}
\end{tikzpicture}
\end{center}

Problem 3 is the shortest in points but its Part B is almost free once Part A is done: every
question in it says ``deduce''. Do not leave it for the last ten minutes.

\section{Making the first sitting feel like the second}

You tend to do better the second time you face an exam than the first. That is not about
knowing more the second time; it is that nothing on the paper is a surprise any more. You can
get most of that effect before the real exam:

\begin{enumerate}
  \item \textbf{Sit the mock once, timed, from blank paper,} before reading the answers here. Mark
  it yourself with this booklet.
  \item \textbf{Read the booklet,} one problem per evening. Read the ``lesson'' boxes twice.
  \item \textbf{Two days later, sit it again} from blank paper, without the booklet. This second
  sitting is the one that counts: it is your resit, done before the exam.
  \item \textbf{Learn the opening routine of each problem type} (the ``How to start'' boxes). In the
  real exam the first line you write is the one you rehearsed, and the rest follows from it.
\end{enumerate}

\begin{remember}[title={When you blank in the room}]
Write the definition. For a differentiability question, write the definition of differentiability.
For a tangent plane, write the formula. For a radial function, write $r=\sqrt{x^2+y^2}$ and
$r_x=x/r$. Definitions are worth marks on their own, and writing one almost always shows you the
next line.
\end{remember}
